%==============================================================================
% DCC062 - Sistemas Operacionais - Apresentacao (Parte Teorica)
% Tema 1: Android e Fuchsia - Grupo 8
%
% Slides enxutos (pouco texto) para a apresentacao oral.
% 1 a 2 slides por ponto do roteiro (12 pontos).
%
% Compilar:  cd apresentacao && latexmk -pdf slides.tex
%   ou:      pdflatex slides.tex   (rodar 2x para o sumario/rodape)
%==============================================================================
\documentclass[aspectratio=169,11pt]{beamer}

%------------------------------------------------------------------ Idioma
\usepackage[utf8]{inputenc}
\usepackage[T1]{fontenc}
\usepackage[brazilian]{babel}
\usepackage{lmodern}

%------------------------------------------------------------------ Visual simples
\usetheme{default}
\usecolortheme{default}
\setbeamertemplate{navigation symbols}{}   % sem os botoes de navegacao
\useinnertheme{circles}                     % marcadores discretos

% Paleta (mesmo azul do trabalho escrito)
\definecolor{azulUFJF}{HTML}{1A5276}
\definecolor{cinzaClaro}{HTML}{F2F4F4}
\definecolor{androidVerde}{HTML}{2E7D32}
\definecolor{fuchsiaCor}{HTML}{8E24AA}

\setbeamercolor{frametitle}{fg=white,bg=azulUFJF}
\setbeamercolor{title}{fg=azulUFJF}
\setbeamercolor{structure}{fg=azulUFJF}
\setbeamercolor{section in toc}{fg=azulUFJF}
\setbeamerfont{frametitle}{series=\bfseries}
\setbeamertemplate{frametitle}{%
  \nointerlineskip
  \begin{beamercolorbox}[wd=\paperwidth,ht=2.6ex,dp=1.1ex,leftskip=0.4cm]{frametitle}
    \insertframetitle
  \end{beamercolorbox}%
}

% Rodape minimalista: tema | secao atual | numero do slide
\setbeamertemplate{footline}{%
  \hbox{%
    \begin{beamercolorbox}[wd=.5\paperwidth,ht=2.4ex,dp=1ex,leftskip=.3cm]{author in head/foot}%
      \tiny Android e Fuchsia --- Grupo 8%
    \end{beamercolorbox}%
    \begin{beamercolorbox}[wd=.5\paperwidth,ht=2.4ex,dp=1ex,rightskip=.3cm]{author in head/foot}%
      \hfill\tiny \insertsectionhead\quad\insertframenumber/\inserttotalframenumber%
    \end{beamercolorbox}%
  }%
}

%------------------------------------------------------- Atalhos para comparacao
% Rotulos coloridos de Android / Fuchsia usados nas colunas
\newcommand{\ANDROID}{\textbf{\color{androidVerde}Android}}
\newcommand{\FUCHSIA}{\textbf{\color{fuchsiaCor}Fuchsia}}

% Bloco de duas colunas Android | Fuchsia.
% Uso: \comparar{itens android}{itens fuchsia}
\newcommand{\comparar}[2]{%
  \begin{columns}[T,onlytextwidth]
    \begin{column}{0.48\textwidth}
      \ANDROID\par\smallskip
      \begin{itemize}\footnotesize #1 \end{itemize}
    \end{column}
    \begin{column}{0.48\textwidth}
      \FUCHSIA\par\smallskip
      \begin{itemize}\footnotesize #2 \end{itemize}
    \end{column}
  \end{columns}%
}

%------------------------------------------------------------------ Metadados
\title{Android e Fuchsia}
\subtitle{Dois sistemas da Google, duas filosofias}
\author[Grupo 8]{%
  \texorpdfstring{%
    Grupo 8 \\ \small Felipe Lazzarini Cunha \\
    Lucas Vieira Resende \\
    Nathan Ferreira da Silva Zoffoli \\
    João Pedro Costa Lopes%
  }{%
    Grupo 8: Felipe Lazzarini Cunha, Lucas Vieira Resende,
    Nathan Ferreira da Silva Zoffoli e João Pedro Costa Lopes%
  }%
}
\institute{\small UFJF --- DCC062 Sistemas Operacionais --- 2026.1}
\date{Julho de 2026}

\begin{document}

%==============================================================================
% CAPA
%==============================================================================
{\setbeamertemplate{footline}{}
\begin{frame}[plain]
  \titlepage
  \begin{center}
    \footnotesize\color{gray}
    \ANDROID: kernel Linux monolítico \quad$\bullet$\quad
    \FUCHSIA: microkernel Zircon
  \end{center}
\end{frame}}

%------------------------------------------------------------------ Roteiro
\begin{frame}{Roteiro}
  \footnotesize
  \begin{columns}[T,onlytextwidth]
    \begin{column}{0.5\textwidth}
      \begin{enumerate}[1.]
        \item Apresentação do tema
        \item Propósito, evolução e virtualização
        \item Tipos de aplicação e uso do kernel
        \item Requisitos das aplicações
        \item Como os requisitos são atendidos
        \item Sistema de interrupções
      \end{enumerate}
    \end{column}
    \begin{column}{0.5\textwidth}
      \begin{enumerate}[1.]\setcounter{enumi}{6}
        \item Suporte a threads
        \item Exclusão mútua, sincronização e IPC
        \item Processos e escalonadores
        \item Memória e memória virtual
        \item I/O e sistemas de arquivos
        \item Características distintivas
      \end{enumerate}
    \end{column}
  \end{columns}
\end{frame}

%==============================================================================
% PONTO 1 - APRESENTACAO
%==============================================================================
\section{1. Apresentação}

\begin{frame}{1. Apresentação do tema}
  Dois sistemas operacionais da \textbf{Google} com filosofias opostas.
  \vfill
  \comparar{
    \item SO móvel mais usado do mundo
    \item \emph{Kernel} Linux modificado (ACK)
    \item Aberto via AOSP
    \item Estreia em 2008 (HTC Dream)
  }{
    \item Construído ``do zero''
    \item Microkernel próprio: \textbf{Zircon}
    \item Segurança por \emph{capacidades}
    \item Público em 2016; Nest Hub em 2021
  }
  \vfill
  \begin{center}\small
    Monolítico e maduro \quad$\times$\quad micronúcleo e modular.
  \end{center}
\end{frame}

%==============================================================================
% PONTO 2 - PROPOSITO, EVOLUCAO, VIRTUALIZACAO
%==============================================================================
\section{2. Propósito}

\begin{frame}{2. Propósito de cada sistema}
  \comparar{
    \item Plataforma completa para dispositivos móveis
    \item \emph{Framework} + \emph{runtime} + serviços
    \item Roda sobre \emph{kernel} Linux adaptado
    \item Foco: ecossistema amplo de apps
  }{
    \item SO de propósito geral, experimental
    \item Modular, atualizável e seguro na base
    \item Tudo é \emph{componente} com capacidades
    \item Foco: isolamento e composição
  }
\end{frame}

\begin{frame}{2. Evolução e virtualização}
  \comparar{
    \item Dalvik $\rightarrow$ \textbf{ART} (\emph{bytecode} DEX)
    \item Modularização: Treble e Mainline
    \item \textbf{AVF/pKVM}: VMs protegidas (Microdroid)
    \item Isola memória mesmo se o host cair
  }{
    \item Magenta $\rightarrow$ \textbf{Zircon} (2017)
    \item Componentes, pacotes e FIDL
    \item \textbf{Starnix}: compatibilidade com Linux
    \item Isolamento por capacidades, não VM
  }
\end{frame}

%==============================================================================
% PONTO 3 - TIPOS DE APLICACAO / USO DO KERNEL   [PREENCHIDO]
%==============================================================================
\section{3. Tipos de aplicação}

\begin{frame}{3. Tipos de aplicação e uso do kernel}
  \comparar{
    \item Apps \emph{interativos} de terceiros
    \item Kernel Linux usado em:
    \item \quad \emph{smartphones} e \emph{tablets}
    \item \quad Android TV, Wear OS
    \item \quad Android Automotive (carros)
    \item Mesmo kernel, muitas classes de aparelho
  }{
    \item Dispositivos conectados / IoT
    \item Zircon usado em:
    \item \quad \emph{smart displays} (Nest Hub)
    \item \quad produtos embarcados da Google
    \item Presença comercial mais restrita
    \item Um só núcleo, foco em segurança
  }
  \vfill
  {\scriptsize\color{gray} Android: apps móveis interativos em muitas formas de dispositivo; Fuchsia: dispositivos conectados, com o Zircon como microkernel comum.}
\end{frame}

%==============================================================================
% PONTO 4 - REQUISITOS
%==============================================================================
\section{4. Requisitos}

\begin{frame}{4. Requisitos das aplicações}
  \comparar{
    \item Baixo \textbf{tempo de resposta} (16,7 ms/quadro)
    \item \textbf{Eficiência energética} (bateria)
    \item Operar sob \textbf{pouca memória}
    \item Isolamento por aplicativo
    \item Portabilidade entre hardwares
  }{
    \item \textbf{Segurança} por capacidades
    \item \textbf{Atualizabilidade} sob demanda
    \item Modularidade e composição
    \item Escalabilidade entre dispositivos
    \item Compatibilidade com Linux
  }
  \vfill
  \begin{center}\small
    Android prioriza o \emph{usuário final}; Fuchsia, a \emph{estrutura} do sistema.
  \end{center}
\end{frame}

%==============================================================================
% PONTO 5 - ATENDIMENTO DOS REQUISITOS
%==============================================================================
\section{5. Atendimento}

\begin{frame}{5. Como o Android atende os requisitos}
  \ANDROID
  \begin{itemize}\footnotesize
    \item \textbf{Responsividade:} Looper/Handler, InputReader/Dispatcher, Choreographer + VSync; alerta \emph{ANR}
    \item \textbf{Energia:} Doze, SystemSuspend, \emph{wakelocks}, mitigação térmica
    \item \textbf{Memória:} \texttt{kswapd} + \texttt{lmkd} encerram processos menos importantes
    \item \textbf{Segurança:} \emph{sandbox} por UID, SELinux, \emph{Verified Boot}, Binder
    \item \textbf{Portabilidade:} HALs e GKI separam plataforma de \emph{drivers}
  \end{itemize}
\end{frame}

\begin{frame}{5. Como o Fuchsia atende os requisitos}
  \FUCHSIA
  \begin{itemize}\footnotesize
    \item \textbf{Segurança:} processos sem autoridade; recursos via \emph{handles} (menor privilégio)
    \item \textbf{Atualizabilidade:} software em pacotes independentes, sob demanda
    \item \textbf{Modularidade:} componentes com \emph{manifesto} e capacidades declaradas
    \item \textbf{Escalabilidade:} \emph{build} configurado por produto
    \item \textbf{Compatibilidade:} Starnix implementa a interface de \emph{syscalls} do Linux
  \end{itemize}
\end{frame}

%==============================================================================
% PONTO 6 - INTERRUPCOES
%==============================================================================
\section{6. Interrupções}

\begin{frame}{6. Sistema de interrupções}
  \comparar{
    \item IRQ $\rightarrow$ \textbf{GIC} $\rightarrow$ CPU
    \item Dividido em duas metades:
    \item \quad \emph{Top Half}: rápido, reconhece a IRQ
    \item \quad \emph{Bottom Half}: SoftIRQ, Tasklet, Workqueue
    \item \emph{Wakeup IRQs} + \emph{wakelocks} p/ acordar
  }{
    \item \emph{Drivers} no espaço de usuário
    \item Zircon só mascara e sinaliza o evento
    \item \emph{Interrupt Object} + \emph{Handle}
    \item Ciclo: \texttt{zx\_interrupt\_wait / ack}
    \item Falha no \emph{driver} não derruba o núcleo
  }
  \vfill
  \begin{center}\small
    Linux: falha de \emph{driver} $=$ \emph{kernel panic}. \quad Zircon: só reinicia o componente.
  \end{center}
\end{frame}

%==============================================================================
% PONTO 7 - THREADS   [PREENCHIDO]
%==============================================================================
\section{7. Threads}

\begin{frame}{7. Suporte a threads}
  \comparar{
    \item \emph{Threads} são do \emph{kernel} Linux (modelo \textbf{1:1})
    \item Base em \texttt{pthreads}; \texttt{Thread} de Java/Kotlin
    \item \emph{UI thread} + \emph{worker threads}
    \item \texttt{HandlerThread} e co-rotinas (Kotlin)
    \item Escalonamento fica a cargo do Linux
  }{
    \item \emph{Thread} é um \textbf{objeto do Zircon}
    \item Um processo contém uma ou mais \emph{threads}
    \item Escalonadas diretamente pelo micronúcleo
    \item \texttt{pthreads}/C11 sobre \texttt{zx\_thread}
    \item \emph{Futexes} para sincronização interna
  }
  \vfill
  {\scriptsize\color{gray} Ambos usam \emph{threads} de núcleo (1:1); o Android as herda do Linux, o Fuchsia as expõe como objetos do Zircon.}
\end{frame}

%==============================================================================
% PONTO 8 - EXCLUSAO MUTUA, SINCRONIZACAO, IPC
%==============================================================================
\section{8. Sincronização e IPC}

\begin{frame}{8. Exclusão mútua, sincronização e IPC}
  \comparar{
    \item Isolamento: 1 processo/UID por app
    \item Monitores Java: \texttt{wait/notify}
    \item \emph{UI thread}: Looper + MessageQueue + Handler
    \item \textbf{Binder} IPC + Service Manager + AIDL
  }{
    \item Recursos via \emph{handles} do Zircon
    \item \emph{Channels}: mensagens \emph{e} \emph{handles}
    \item \emph{Sockets}: fluxo de \emph{bytes}
    \item \emph{Futexes} locais ao processo
  }
  \vfill
  \begin{center}\small
    Android: IPC centralizado no Binder. \quad Fuchsia: canais + capacidades.
  \end{center}
\end{frame}

%==============================================================================
% PONTO 9 - PROCESSOS E ESCALONADORES
%==============================================================================
\section{9. Processos}

\begin{frame}{9. Processos e escalonadores}
  \comparar{
    \item Processos criados pelo \textbf{Zygote}
    \item \texttt{fork()} + \emph{copy-on-write}
    \item Escalonamento é do \emph{kernel} Linux
    \item Prioridade: primeiro plano $\rightarrow$ vazio
  }{
    \item Kernel escalona \emph{threads}, não apps
    \item Apps $=$ \emph{componentes} via \emph{runners}
    \item ELF Runner, Web Runner, Starnix
    \item Ciclo de vida no \emph{Component Manager}
  }
  \vfill
  \begin{center}\footnotesize
    \ANDROID: Zygote ``aquecido'' acelera a abertura de apps compartilhando páginas.
  \end{center}
\end{frame}

%==============================================================================
% PONTO 10 - MEMORIA
%==============================================================================
\section{10. Memória}

\begin{frame}{10. Memória e memória virtual}
  \comparar{
    \item Paginação sob demanda + ASLR
    \item \textbf{Sem swap em disco} $\rightarrow$ \textbf{ZRAM} (comprimida)
    \item Métrica \textbf{PSS} mede custo real de RAM
    \item \texttt{lmkd} + \emph{OOM score} liberam memória
  }{
    \item \textbf{VMO}: contêiner de páginas físicas
    \item \textbf{VMAR}: árvore do espaço de endereços
    \item Acesso só com \emph{handle} válido
    \item Falha de \emph{driver} contida na sua VMAR
  }
\end{frame}

%==============================================================================
% PONTO 11 - I/O E SISTEMAS DE ARQUIVOS   [PREENCHIDO]
%==============================================================================
\section{11. I/O e arquivos}

\begin{frame}{11. I/O e sistemas de arquivos}
  \comparar{
    \item I/O pela \textbf{VFS} do Linux
    \item \textbf{F2FS} e ext4 (flash interno)
    \item EROFS (partição de sistema, só leitura)
    \item FAT/exFAT (cartões SD) via FUSE
    \item \emph{Scoped storage} isola arquivos por app
  }{
    \item Sistemas de arquivos em \emph{espaço de usuário}
    \item Acesso via \emph{canais} / FIDL
    \item \textbf{Fxfs}: sistema de arquivos padrão atual
    \item \textbf{Blobfs}: pacotes por \emph{hash} de conteúdo
    \item FVM gerencia os volumes
  }
  \vfill
  {\scriptsize\color{gray} Android reaproveita a VFS do Linux (F2FS/ext4/EROFS); no Fuchsia cada FS é um componente isolado (Fxfs, Blobfs).}
\end{frame}

%==============================================================================
% PONTO 12 - CARACTERISTICAS DISTINTIVAS
%==============================================================================
\section{12. Distintivos}

\begin{frame}{12. Características distintivas}
  \comparar{
    \item \textbf{Zygote}: processo ``aquecido'' com ART
    \item \emph{fork} + \emph{copy-on-write} $\rightarrow$ boot rápido de apps
    \item \textbf{Treble} isola SO do código do fabricante
    \item \textbf{APEX/Mainline}: atualiza módulos pela loja
  }{
    \item Boot cresce como \textbf{árvore de componentes}
    \item Cada app traz seu próprio \emph{runner}
    \item \textbf{Software efêmero}: pacotes por URL
    \item Sistema \emph{evergreen}, sem reinstalar
  }
\end{frame}

%==============================================================================
% ENCERRAMENTO
%==============================================================================
\section{}

\begin{frame}{Síntese}
  \begin{center}
    \large Uma mesma empresa, dois caminhos.
  \end{center}
  \vfill
  \comparar{
    \item Monolítico (Linux), maduro e onipresente
    \item Otimiza a \emph{experiência do usuário}
    \item Evolui adaptando um núcleo consolidado
  }{
    \item Micronúcleo (Zircon), moderno e modular
    \item Otimiza \emph{segurança} e \emph{atualização}
    \item Repensa o SO a partir do zero
  }
  \vfill
  \begin{center}\small\color{azulUFJF}
    Desempenho e simplicidade \quad$\times$\quad isolamento e atualizabilidade.
  \end{center}
\end{frame}

{\setbeamertemplate{footline}{}
\begin{frame}[plain]
  \centering
  \vfill
  {\Huge\color{azulUFJF}\bfseries Obrigado!}\\[2em]
  \footnotesize Grupo 8 --- DCC062 Sistemas Operacionais --- UFJF --- 2026.1
  \vfill
\end{frame}}

\end{document}
